\DocumentMetadata{
  pdfversion=2.0,
  pdfstandard=ua-2,
  lang=en-US,
  tagging=on,
  tagging-setup={math/setup=mathml-SE,tabsorder=structure}
}
\documentclass[11pt,letterpaper]{article}
\usepackage[margin=0.68in,includefoot]{geometry}
\usepackage{newcomputermodern}
\usepackage{amsmath,amscd,mathtools}
\providecommand{\square}{\mdlgwhtsquare}
\usepackage[hidelinks]{hyperref}
\hypersetup{
  pdftitle={Analysis First-Year Exam, August 2016, Part 2},
  pdfauthor={University of Florida Department of Mathematics},
  pdfsubject={Graduate mathematics examination},
  pdfdisplaydoctitle=true,
  bookmarksopen=true
}
\setcounter{secnumdepth}{0}
\setlength{\parindent}{0pt}
\setlength{\parskip}{0.28em}
\setlength{\emergencystretch}{3em}
\setlength{\leftmargini}{2.15em}
\setlength{\leftmarginii}{2.65em}
\setlength{\leftmarginiii}{3em}
\renewcommand{\labelenumi}{\textbf{\theenumi.}}
\pagestyle{empty}
\raggedbottom
\allowdisplaybreaks
\newenvironment{examproblems}[1]{%
  \begin{enumerate}%
  \setcounter{enumi}{#1}%
  \setlength{\topsep}{0.35em}%
  \setlength{\partopsep}{0pt}%
  \setlength{\itemsep}{0.48em}%
  \setlength{\parsep}{0.18em}%
}{\end{enumerate}}
\newenvironment{examsubparts}{%
  \begin{enumerate}%
  \setlength{\topsep}{0.18em}%
  \setlength{\partopsep}{0pt}%
  \setlength{\itemsep}{0.16em}%
  \setlength{\parsep}{0.1em}%
}{\end{enumerate}}
\newenvironment{examinstructions}{%
  \par\begingroup\itshape
}{%
  \par\endgroup\vspace{0.18em}
}
\makeatletter
\renewcommand\section{\@startsection{section}{1}{0pt}%
  {-1.25ex plus -.25ex minus -.1ex}%
  {0.55ex plus .1ex}%
  {\normalfont\large\bfseries}}
\renewcommand{\@maketitle}{%
  \newpage\null\vskip -1.2em%
  {\footnotesize\bfseries
    \href{https://www.ufl.edu/}{University of Florida}, \href{https://math.ufl.edu/}{Department of Mathematics}\par}%
  \vskip 0.18em%
  \tagstructbegin{tag=Title}%
    {\LARGE\bfseries\href{https://gma.math.ufl.edu/exams/analysis-first-year/}{Analysis First-Year Exam}, August 2016, Part 2\par}%
  \tagstructend%
  \vskip 0.35em%
  \tagmcbegin{artifact}%
    \hrule height 1pt%
  \tagmcend%
  \vskip 0.55em%
}
\makeatother
\title{Analysis First-Year Exam, August 2016, Part 2}
\author{}
\date{}
\begin{document}
\maketitle
\thispagestyle{empty}
\begin{examinstructions}
Answer exactly FOUR questions. Write solutions in a neat and logical fashion, giving complete reasons for all steps.
\end{examinstructions}
\begin{examproblems}{0}
\item
Let the function \(f:\mathbb{R}\to\mathbb{R}\) be uniformly continuous. For each positive integer~\(n\) let \(f_n(t)=f\!\left(t+\frac{1}{n}\right)\) whenever~\(t\) is a real number. Prove that the sequence \((f_n)_{n=1}^{\infty}\) is uniformly convergent. Give an example (with justification) to show that the conclusion can fail if the hypothesis uniformly is dropped.
\item
Let \(A=\{a_n:n\geq 1\}\) be a countably infinite subset of \([0,1]\) and let \(1_A:[0,1]\to\mathbb{R}\) be its indicator (or characteristic) function. Exhibit such a set~\(A\) for which \(1_A\) is not Riemann integrable and exhibit such a set~\(A\) for which \(1_A\) is Riemann integrable, providing justification in each case.
\item
Let \((f_n)_{n=1}^{\infty}\) be a sequence of continuous functions from \([0,1]\) to \([0,1]\). Prove that if \(f_n(t)\) decreases to~\(0\) whenever \(0\leq t\leq 1\), then \(\int_0^1 f_n(t)\,dt\to 0\). Does the same conclusion follow when decreases is replaced by converges? Justify.
\item
Prove that if the function \(f:\mathbb{R}\to\mathbb{R}\) is differentiable then its derivative is measurable.
\item
Let~\(f\) be a Lebesgue integrable function on~\(\mathbb{R}\) and let~\(g\) be defined on~\(\mathbb{R}\) by \(g(y)=\int_{-\infty}^{\infty}\cos(xy)f(x)\,dx\). Prove that \(\lim_{k\to\infty}g(k)=0\).
\end{examproblems}
\end{document}
